% ============================================================
\chapter{Th\'eor\`eme de Radon-Nikodym}
\label{ch:radon_nikodym}
% ============================================================

\section{Introduction}

Quand une mesure $\nu$ est-elle une ``version pond\'er\'ee'' d'une mesure $\mu$~? C'est-\`a-dire, quand existe-t-il une fonction $f$ telle que $\nu(A) = \int_A f\, d\mu$ pour tout ensemble mesurable $A$~? Johann Radon (1913) et Otton Nikodym (1930) r\'epondent~: si et seulement si $\nu$ est absolument continue par rapport \`a $\mu$, ce qui signifie que $\mu(A) = 0$ entra\^ine $\nu(A) = 0$. La fonction $f$, appel\'ee \emph{d\'eriv\'ee de Radon--Nikodym}, g\'en\'eralise la notion de densit\'e de probabilit\'e. Ce th\'eor\`eme est l'un des piliers de l'analyse moderne~: il fonde la d\'efinition rigoureuse de l'esp\'erance conditionnelle en probabilit\'es, la dualit\'e des espaces $L^p$, et le changement de mesure dans le calcul stochastique (th\'eor\`eme de Girsanov).

\begin{intuition}[title=\textbf{Radon-Nikodym~: changer de r\'ef\'erentiel}]
Si $\nu$ est absolument continue par rapport \`a $\mu$, cela signifie que $\nu$ ne
\og voit \fg que ce que $\mu$ voit~: les ensembles n\'egligeables pour $\mu$ le sont
aussi pour $\nu$. Le th\'eor\`eme de Radon-Nikodym affirme alors que $\nu$ s'obtient
en \og pesant \fg la mesure $\mu$ par une densit\'e~: $d\nu = f\, d\mu$.
C'est la g\'en\'eralisation du fait que toute probabilit\'e sur $\R$ \`a densit\'e
s'\'ecrit $\PP(A) = \int_A f(x)\, dx$.
\end{intuition}

% ------------------------------------------------------------
\section{Continuit\'e absolue}
% ------------------------------------------------------------

\begin{definition}[Continuit\'e absolue]
\label{def:abs_cont}
\index{continuit\'e absolue}
Soient $\mu$ une mesure positive et $\nu$ une mesure sign\'ee sur $(X, \mathcal{A})$.
On dit que $\nu$ est \emph{absolument continue} par rapport \`a $\mu$, not\'e
$\nu \ll \mu$\index{$\ll$}, si~:
\[
    \forall A \in \mathcal{A}, \quad \mu(A) = 0 \implies \nu(A) = 0.
\]
\end{definition}

\begin{proposition}[Caract\'erisation $\varepsilon$-$\delta$]
\label{prop:eps_delta}
Si $\mu$ est une mesure positive et $\nu$ est une mesure sign\'ee \emph{finie},
alors $\nu \ll \mu$ si et seulement si~:
\[
    \forall \varepsilon > 0,\, \exists\, \delta > 0,\, \forall A \in \mathcal{A},\quad
    \mu(A) < \delta \implies \abs{\nu(A)} < \varepsilon.
\]
\end{proposition}

\begin{proof}
($\Leftarrow$) \'Evident~: si $\mu(A) = 0$, alors $\mu(A) < \delta$ pour tout
$\delta$, d'o\`u $\abs{\nu(A)} < \varepsilon$ pour tout $\varepsilon$, soit $\nu(A) = 0$.

($\Rightarrow$) Par l'absurde. Supposons que $\nu \ll \mu$ mais que la condition
$\varepsilon$-$\delta$ \'echoue. Il existe $\varepsilon_0 > 0$ et une suite $(A_n)$
avec $\mu(A_n) < 2^{-n}$ et $\abs{\nu(A_n)} \geq \varepsilon_0$. Posons
$B_n = \bigcup_{k \geq n} A_k$. Alors $\mu(B_n) \leq \sum_{k \geq n} 2^{-k} = 2^{1-n}$
et $B = \bigcap_n B_n$ v\'erifie $\mu(B) = 0$. Par continuit\'e d\'ecroissante,
$\abs{\nu}(B) = \lim_n \abs{\nu}(B_n)$. Or $\abs{\nu}(B_n) \geq \abs{\nu(A_n)}
\geq \varepsilon_0$ pour tout $n$, donc $\abs{\nu}(B) \geq \varepsilon_0 > 0$,
contredisant $\nu \ll \mu$.
\end{proof}

\begin{exemple}
\begin{enumerate}
    \item Si $f \geq 0$ est mesurable et $\nu(A) = \int_A f\, d\mu$, alors $\nu \ll \mu$.
    \item La mesure de Dirac $\delta_0$ n'est pas absolument continue par rapport \`a
    la mesure de Lebesgue $\lambda$ sur $\R$~: $\lambda(\{0\}) = 0$ mais $\delta_0(\{0\}) = 1$.
    \item Si $\nu \ll \mu$ et $\mu \ll \nu$, on dit que $\mu$ et $\nu$ sont
    \emph{\'equivalentes}, not\'e $\mu \sim \nu$.
\end{enumerate}
\end{exemple}

% ------------------------------------------------------------
\section{Th\'eor\`eme de Radon-Nikodym}
% ------------------------------------------------------------

\begin{theoreme}[Radon-Nikodym]
\label{thm:radon_nikodym}
\index{th\'eor\`eme de Radon-Nikodym}
Soit $(X, \mathcal{A}, \mu)$ un espace mesur\'e $\sigma$-fini et soit $\nu$ une
mesure sign\'ee $\sigma$-finie sur $(X, \mathcal{A})$ avec $\nu \ll \mu$. Alors il
existe une fonction mesurable $f : X \to [-\infty, +\infty]$, int\'egrable par rapport
\`a $\mu$ (si $\nu$ est finie), telle que~:
\[
    \nu(A) = \int_A f\, d\mu, \qquad \forall A \in \mathcal{A}.
\]
De plus, $f$ est unique $\mu$-presque partout. Si $\nu$ est une mesure positive,
alors $f \geq 0$ $\mu$-p.p.
\end{theoreme}

\begin{proof}[Preuve (cas $\mu, \nu$ mesures positives finies)]
On pr\'esente la preuve de Von Neumann.

Posons $\lambda = \mu + \nu$. C'est une mesure positive finie. Pour tout
$g \in L^2(\lambda)$, la forme lin\'eaire
$\varphi(g) = \int g\, d\nu$ est continue sur $L^2(\lambda)$ car, par
Cauchy--Schwarz~:
\[
    \abs{\varphi(g)} = \left|\int g\, d\nu\right| \leq \int \abs{g}\, d\nu
    \leq \int \abs{g}\, d\lambda
    \leq \left(\int g^2\, d\lambda\right)^{1/2} \lambda(X)^{1/2}.
\]
Par le th\'eor\`eme de repr\'esentation de Riesz pour les espaces de Hilbert, il
existe $h \in L^2(\lambda)$ tel que~:
\[
    \int g\, d\nu = \int gh\, d\lambda, \qquad \forall g \in L^2(\lambda).
\]
En prenant $g = \ind_A$, on obtient $\nu(A) = \int_A h\, d\lambda$ pour tout
$A \in \mathcal{A}$. Comme $\nu(A) \geq 0$ et $\nu(A) \leq \lambda(A)$ pour tout $A$,
on d\'eduit que $0 \leq h \leq 1$ $\lambda$-p.p.

On peut \'ecrire~: $\nu(A) = \int_A h\, d\mu + \int_A h\, d\nu$, soit
$\int_A (1-h)\, d\nu = \int_A h\, d\mu$. En posant
$E_0 = \{h = 1\}$ et en v\'erifiant que $\mu(E_0) = 0$ (car
$0 = \int_{E_0}(1-h)\, d\nu = \int_{E_0} h\, d\mu = \mu(E_0)$), on d\'efinit
$f = h/(1-h)$ sur $E_0^c$ et $f = 0$ sur $E_0$. Le th\'eor\`eme de convergence
monotone appliqu\'e \`a la suite $g_n = (1 + h + \cdots + h^n) \ind_A$ donne~:
\[
    \nu(A) = \int_A f\, d\mu.
\]
L'unicit\'e de $f$ $\mu$-p.p.\ r\'esulte du fait que si $\int_A f\, d\mu = \int_A g\, d\mu$
pour tout $A$, alors $f = g$ $\mu$-p.p.
\end{proof}

\begin{definition}[D\'eriv\'ee de Radon-Nikodym]
\index{d\'eriv\'ee de Radon-Nikodym}
La fonction $f$ du th\'eor\`eme~\ref{thm:radon_nikodym} est appel\'ee la
\emph{d\'eriv\'ee de Radon-Nikodym} de $\nu$ par rapport \`a $\mu$ et est not\'ee~:
\[
    f = \frac{d\nu}{d\mu}.
\]
\end{definition}

\begin{proposition}[Propri\'et\'es de la d\'eriv\'ee de Radon-Nikodym]
\label{prop:rn_props}
Soient $\mu, \nu, \lambda$ des mesures $\sigma$-finies.
\begin{enumerate}
    \item \textbf{R\`egle de la cha\^ine~:} si $\nu \ll \mu \ll \lambda$, alors
    $\dfrac{d\nu}{d\lambda} = \dfrac{d\nu}{d\mu} \cdot \dfrac{d\mu}{d\lambda}$
    $\lambda$-p.p.
    \item \textbf{Inverse~:} si $\nu \ll \mu$ et $\mu \ll \nu$, alors
    $\dfrac{d\mu}{d\nu} = \left(\dfrac{d\nu}{d\mu}\right)^{-1}$ $\mu$-p.p.
    \item \textbf{Lin\'earit\'e~:} si $\nu_1, \nu_2 \ll \mu$, alors
    $\dfrac{d(\alpha\nu_1 + \beta\nu_2)}{d\mu} = \alpha\dfrac{d\nu_1}{d\mu}
    + \beta\dfrac{d\nu_2}{d\mu}$.
\end{enumerate}
\end{proposition}

\begin{formules}[title=\textbf{Formules cl\'es de Radon-Nikodym}]
Si $\nu \ll \mu$ avec $f = \frac{d\nu}{d\mu}$, alors pour toute fonction
$g$ mesurable positive ou $\nu$-int\'egrable~:
\[
    \int g\, d\nu = \int g \cdot f\, d\mu = \int g \cdot \frac{d\nu}{d\mu}\, d\mu.
\]
R\`egle de la cha\^ine~: $\dfrac{d\nu}{d\lambda} = \dfrac{d\nu}{d\mu} \cdot \dfrac{d\mu}{d\lambda}$.
\end{formules}

% ------------------------------------------------------------
\section{Th\'eor\`eme de d\'ecomposition de Lebesgue}
% ------------------------------------------------------------

\begin{theoreme}[D\'ecomposition de Lebesgue]
\label{thm:lebesgue_decomp}
\index{d\'ecomposition de Lebesgue}
Soit $(X, \mathcal{A}, \mu)$ un espace mesur\'e $\sigma$-fini et soit $\nu$ une mesure
sign\'ee $\sigma$-finie. Alors il existe un unique couple de mesures sign\'ees
$(\nu_a, \nu_s)$ tel que~:
\[
    \nu = \nu_a + \nu_s, \qquad \nu_a \ll \mu, \qquad \nu_s \perp \mu.
\]
De plus, $\nu_a(A) = \int_A f\, d\mu$ pour une certaine $f$ mesurable (la d\'eriv\'ee
de Radon-Nikodym de $\nu_a$ par rapport \`a $\mu$).
\end{theoreme}

\begin{proof}[Esquisse de preuve]
Il suffit de traiter le cas o\`u $\nu$ est une mesure positive finie et $\mu$ est finie.
Pour chaque $\lambda > 0$, la mesure sign\'ee $\nu - \lambda\mu$ admet une
d\'ecomposition de Hahn~: $X = P_\lambda \cup N_\lambda$ avec $P_\lambda$ positif et
$N_\lambda$ n\'egatif pour $\nu - \lambda\mu$. On peut choisir $P_\lambda$ croissant
en $\lambda$ (quitte \`a remplacer $P_\lambda$ par $\bigcup_{\lambda' \leq \lambda} P_{\lambda'}$).
Posons $\lambda_0 = \sup\{\lambda > 0 : \nu(N_\lambda) > 0\}$ et consid\'erons une suite
$\lambda_n \nearrow \lambda_0$. L'ensemble $N = \bigcap_{n} N_{\lambda_n}$ satisfait
$\mu(N) = 0$~: en effet, pour tout $n$, $\nu(A) \leq \lambda_n \mu(A)$ pour tout
$A \subseteq P_{\lambda_n}^c$, et en passant \`a la limite on obtient $\mu(N) = 0$.
On pose $\nu_s(\cdot) = \nu(\cdot \cap N)$ et $\nu_a(\cdot) = \nu(\cdot \cap N^c)$.
Comme $\mu(N) = 0$, on a $\nu_s \perp \mu$. Pour montrer que $\nu_a \ll \mu$~: si
$\mu(A) = 0$ avec $A \subseteq N^c$, alors $A \subseteq P_{\lambda_n}$ pour $n$ assez
grand, et $\nu(A) \leq \lambda_n \mu(A) = 0$. Le th\'eor\`eme de Radon-Nikodym
appliqu\'e \`a $\nu_a$ fournit la densit\'e $f$.
Voir~\textsc{Rudin}, \emph{Real and Complex Analysis}, Theorem~6.10.
\end{proof}

% ------------------------------------------------------------
\section{Applications en probabilit\'es}
% ------------------------------------------------------------

\begin{exemple}[Densit\'e d'une variable al\'eatoire]
Soit $X$ une variable al\'eatoire r\'eelle sur $(\Omega, \mathcal{F}, \PP)$ dont la
loi $\PP_X$ est absolument continue par rapport \`a la mesure de Lebesgue $\lambda$.
Alors il existe $f_X \geq 0$ avec $\int f_X\, d\lambda = 1$ telle que~:
\[
    \PP(X \in A) = \int_A f_X(x)\, dx, \qquad \forall A \in \mathcal{B}(\R).
\]
La fonction $f_X = \frac{d\PP_X}{d\lambda}$ est la \emph{densit\'e} de $X$.
\end{exemple}

\begin{exemple}[Esp\'erance conditionnelle]
Soit $(\Omega, \mathcal{F}, \PP)$ un espace probabilis\'e, $\mathcal{G} \subset \mathcal{F}$
une sous-tribu, et $Y \in L^1(\PP)$. La mesure sign\'ee
$\nu(A) = \int_A Y\, d\PP$, $A \in \mathcal{G}$, est absolument continue par rapport
\`a $\PP|_{\mathcal{G}}$. Par Radon-Nikodym, il existe une unique (p.s.) fonction
$\mathcal{G}$-mesurable $Z \in L^1(\PP)$ telle que~:
\[
    \int_A Z\, d\PP = \int_A Y\, d\PP, \qquad \forall A \in \mathcal{G}.
\]
Cette fonction $Z$ est l'\emph{esp\'erance conditionnelle} $\E[Y | \mathcal{G}]$.
\end{exemple}

\begin{exemple}[Rapport de vraisemblance]
Soient $\PP$ et $\Q$ deux probabilit\'es sur $(\Omega, \mathcal{F})$ avec $\Q \ll \PP$.
Le \emph{rapport de vraisemblance} est $L = \frac{d\Q}{d\PP}$, et pour toute variable
$Y$ $\Q$-int\'egrable~:
\[
    \E_{\Q}[Y] = \E_{\PP}[Y \cdot L].
\]
C'est la base du changement de mesure en finance (th\'eor\`eme de Girsanov) et en
statistique (test de Neyman--Pearson).
\end{exemple}

% ------------------------------------------------------------
\section{Exercices}
% ------------------------------------------------------------

\begin{exercice}
Soit $\mu$ la mesure de Lebesgue sur $[0,1]$ et $\nu(A) = \int_A 2x\, dx$.
V\'erifier que $\nu \ll \mu$, identifier $\frac{d\nu}{d\mu}$, et calculer
$\int_{[0,1]} x^2\, d\nu$.
\end{exercice}

\begin{exercice}
Montrer que si $\nu \ll \mu$ et $\nu \perp \mu$, alors $\nu = 0$.
\end{exercice}

\begin{exercice}
Soient $\mu$ et $\nu$ deux mesures positives $\sigma$-finies avec $\nu \ll \mu$.
Montrer que $\abs{\nu}(A) = \int_A \abs{\frac{d\nu}{d\mu}}\, d\mu$.
\end{exercice}

\begin{exercice}[R\`egle de la cha\^ine]
Soient $\lambda$ la mesure de Lebesgue, $\mu(A) = \int_A e^{-x} \ind_{[0,\infty)}(x)\, d\lambda(x)$,
et $\nu(A) = \int_A x e^{-x} \ind_{[0,\infty)}(x)\, d\lambda(x)$.
Calculer $\frac{d\nu}{d\mu}$ et $\frac{d\mu}{d\lambda}$.
V\'erifier la r\`egle de la cha\^ine.
\end{exercice}

\begin{exercice}
Soit $(\Omega, \mathcal{F}, \PP)$ un espace probabilis\'e et $X$ une variable al\'eatoire
avec $\E[\abs{X}] < \infty$. Soit $\mathcal{G} = \sigma(Y)$ o\`u $Y$ est une variable
al\'eatoire discr\`ete prenant les valeurs $y_1, y_2, \ldots$ Montrer que~:
\[
    \E[X | \mathcal{G}](\omega) = \sum_{k} \frac{\E[X \ind_{\{Y = y_k\}}]}{\PP(Y = y_k)}
    \ind_{\{Y = y_k\}}(\omega).
\]
\end{exercice}

\begin{exercice}
Montrer l'unicit\'e de la d\'ecomposition de Lebesgue~: si $\nu = \nu_a + \nu_s$
et $\nu = \nu_a' + \nu_s'$ avec $\nu_a, \nu_a' \ll \mu$ et
$\nu_s, \nu_s' \perp \mu$, alors $\nu_a = \nu_a'$ et $\nu_s = \nu_s'$.
\end{exercice}
